\documentclass[aspectratio=169,10pt,english]{beamer}

\input{../../../_preambulo_beamer.tex}
\input{../../../_comandos_beamer.tex}

\selectlanguage{english}

\title{MA1001B - Statistical Modeling for Decision Making}
\subtitle{Lesson 20: Continuous Uniform Distribution and Monte Carlo}
\author{}
\institute{Tecnol\'ogico de Monterrey}
\date{}

\begin{document}

\begin{frame}[plain]
  \vspace{1.2cm}
  {\Large\bfseries MA1001B - Statistical Modeling for Decision Making\par}
  \vspace{0.7cm}
  {\large Lesson 20: Continuous Uniform and Monte Carlo\par}
  \vspace{0.7cm}
  {\color{TecRojo}\rule{\textwidth}{1.2pt}}
  \vspace{0.8cm}

  MA1001B Faculty\\[0.3cm]
  Term\\[0.3cm]
  Tecnol\'ogico de Monterrey
\end{frame}

\begin{frame}{The Flat-Density Question}
  \begin{alertblock}{Guiding question}
    If processing time is equally likely anywhere on $[8,20]$ minutes, how
    is $P(X>16)$ computed, and how can simulation check it?
  \end{alertblock}

  \vspace{0.6em}
  By the end of this lesson, you can:
  \begin{itemize}
    \item write the flat density $f(x)=1/(b-a)$;
    \item compute a tail as remaining length over $b-a$;
    \item check that tail with a seeded Monte Carlo sample;
    \item explain why a piled density breaks the uniform model.
  \end{itemize}

  \vfill
  \scriptsize All processing times used today are synthetic. No real company data are used.
\end{frame}

\begin{frame}{Uniform$(8,20)$ Processing Minutes}
  \small
  Density height and mean:
  \[
    f(x)=\frac{1}{20-8}=\frac{1}{12}=0.083333,
    \qquad
    E[X]=\frac{8+20}{2}=14.
  \]
  Tail by length:
  \[
    P(X>16)=\frac{20-16}{12}=\frac{1}{3}=0.333333.
  \]

  \vspace{0.3em}
  \begin{exampleblock}{Synthetic ticket clock}
    Every minute between 8 and 20 is equally likely. Probability is length,
    not a point height.
  \end{exampleblock}
\end{frame}

\begin{frame}[fragile]{Step 1: Exact Uniform Tail}
  \begin{columns}[T]
    \begin{column}{0.42\textwidth}
      \small
      \begin{lstlisting}
model = uniform(
  loc=8, scale=12
)
p_gt_16 = 1 - model.cdf(16)
      \end{lstlisting}

      \begin{block}{Verified output}
        Height $=0.083333$\\
        Mean $=14.000000$\\
        Length ratio $=0.333333$\\
        $P(X>16)=0.333333$
      \end{block}
    \end{column}
    \begin{column}{0.58\textwidth}
      \centering
      \includegraphics[width=\textwidth]{../figures/uniform_mc_01_exact_uniform.png}
      \vspace{0.2em}
      \scriptsize Flat density and shaded tail from Step 1
    \end{column}
  \end{columns}
\end{frame}

\begin{frame}{Monte Carlo Recovers Length Probability}
  \small
  Draw $n=100{,}000$ times from $\mathrm{Uniform}(8,20)$ with seed 42.
  Count the fraction of draws above 16:
  \[
    \widehat{P}(X>16)=\frac{\#\{X_i>16\}}{100{,}000}.
  \]
  The simulated tail is $0.333290$. Absolute error versus $1/3$ is
  $0.000043$. The simulated mean is $14.007499$ versus exact $14$.
\end{frame}

\begin{frame}[fragile]{Step 2: $100{,}000$ Seed-42 Draws}
  \begin{columns}[T]
    \begin{column}{0.40\textwidth}
      \small
      \begin{lstlisting}
draws = model.rvs(
  size=100_000,
  random_state=rng,
)
p_hat = np.mean(draws > 16)
      \end{lstlisting}

      \begin{block}{Verified output}
        Exact tail $=0.333333$\\
        Monte Carlo $=0.333290$\\
        Abs.\ error $=0.000043$\\
        MC mean $=14.007499$
      \end{block}
    \end{column}
    \begin{column}{0.60\textwidth}
      \centering
      \includegraphics[width=\textwidth]{../figures/uniform_mc_02_monte_carlo.png}
      \vspace{0.2em}
      \scriptsize Histogram versus exact density from Step 2
    \end{column}
  \end{columns}
\end{frame}

\begin{frame}{Step 3: Uniform Fails When Density Piles}
  \begin{columns}[T]
    \begin{column}{0.42\textwidth}
      \small
      Uniform $P(X>16)=0.333333$.

      \vspace{0.4em}
      Right-triangular pile at 20:
      \[
        P(X>16)=0.555556.
      \]
      Piled mean $=16$, not $14$.

      \vspace{0.4em}
      \textbf{Limit:} a flat model is wrong when density piles at one end.
      Lesson 21 standardizes a bell-shaped clock.
    \end{column}
    \begin{column}{0.58\textwidth}
      \centering
      \includegraphics[width=\textwidth]{../figures/uniform_mc_03_piled_density_limit.png}
      \vspace{0.2em}
      \scriptsize Flat versus piled densities from Step 3
    \end{column}
  \end{columns}
\end{frame}

\begin{frame}{Concept Check}
  \begin{enumerate}
    \item Why is $f(x)=1/12$ on $[8,20]$ a density, not a probability?
    \item Compute $P(X>16)$ from remaining length $4$ and total length $12$.
    \item The Monte Carlo tail is $0.333290$. Does that replace $1/3$?
    \item Why is $0.555556$ the wrong number for a uniform clock?
  \end{enumerate}

  \vspace{0.6em}
  \begin{block}{Connection to Lesson 21}
    The session can continue directly into the standard normal curve and
    $z$-scores.
  \end{block}

  \vfill
  \scriptsize Source: OpenStax, \textit{Introductory Business Statistics 2e},
  Chapter 5, Section 5.2. CC BY-NC-SA.
\end{frame}

\appendix



\end{document}
